\documentclass[11pt,a4paper]{article}
\usepackage[T1]{fontenc}
\usepackage[utf8]{inputenc}
\usepackage[margin=27mm,headheight=14pt]{geometry}
\usepackage{amsmath,amssymb,amsthm,mathtools}
\usepackage{newtxtext,newtxmath}
\usepackage{microtype}
\usepackage{booktabs,longtable,tabularx,array}
\usepackage{xcolor}
\usepackage[numbers,sort&compress]{natbib}
\usepackage{xurl}
\usepackage{fancyhdr}
\usepackage{enumitem}
\usepackage{hyperref}
\definecolor{linkblue}{RGB}{27,62,93}
\hypersetup{colorlinks=true,linkcolor=linkblue,citecolor=linkblue,urlcolor=linkblue,
 pdftitle={A logarithmic-budget sieve for Erdos problem 647},
 pdfauthor={GPT 6 Astra; Justin V. Wick (operator)},
 pdfsubject={Explicit sparsity bounds and Lean formalization},pdfkeywords={Erdos 647, sieve, divisor function, Lean, formal verification}}
\urlstyle{tt}
\setlength{\emergencystretch}{2em}
\setlength{\parskip}{3pt}
\setlength{\parindent}{1.25em}
\setcounter{tocdepth}{1}
\pagestyle{fancy}
\fancyhf{}
\fancyhead[L]{\small A logarithmic-budget sieve for Erd\H{o}s \#647}
\fancyhead[R]{\small Astra / Wick}
\fancyfoot[C]{\thepage}
\renewcommand{\headrulewidth}{0.3pt}
\newtheorem{theorem}{Theorem}[section]
\newtheorem{lemma}[theorem]{Lemma}
\newtheorem{proposition}[theorem]{Proposition}
\newtheorem{corollary}[theorem]{Corollary}
\theoremstyle{definition}
\newtheorem{definition}[theorem]{Definition}
\theoremstyle{remark}
\newtheorem{remark}[theorem]{Remark}
\numberwithin{equation}{section}
\newcommand{\N}{\mathbb N}
\newcommand{\Z}{\mathbb Z}
\newcommand{\R}{\mathbb R}
\newcommand{\Pset}{\mathcal P}
\newcommand{\Cand}{\mathcal C}
\newcommand{\ind}{\mathbf 1}
\newcommand{\rate}{\kappa}
\newcommand{\smallcode}[1]{{\small\texttt{#1}}}
\newcommand{\leanname}[1]{{\small\nolinkurl{#1}}}
\newcommand{\floor}[1]{\left\lfloor #1\right\rfloor}
\newcommand{\ceil}[1]{\left\lceil #1\right\rceil}
\DeclareMathOperator{\eprod}{E}
\title{\vspace{-1.2em}\textbf{A logarithmic-budget sieve for Erd\H{o}s problem \#647}\\[0.4em]
\large An explicit sparsity bound with a Lean formalization}
\author{\textbf{GPT 6 Astra}\\[-1pt]\small AI author\\[0.5em]
\textbf{Justin V. Wick}\\[-1pt]\small Operator\\[0.5em]
\small Independent project; no institutional affiliation}
\date{September 2026}
\begin{document}
\maketitle
\begin{abstract}
Erd\H{o}s and Selfridge asked whether any integer $n>24$ satisfies
$\tau(n-k)\le k+2$ for every $1\le k<n$, where $\tau$ counts positive divisors.
Let $C(X)$ count integers satisfying this full condition up to $X$.
We give a proof, accompanied by a Lean formalization, that for every fixed
$0<c<1499/10^6$ and all sufficiently large natural $X$,
\[
 C(X)\le X\exp\!\left(-c\,\frac{(\log X)^{\log 2/(1+\log 2)}}{\log\log X}\right).
\]
The argument converts the divisor restrictions into a logarithmic budget for
distinct prime factors, subtracts unavoidable small-prime occurrences, and uses
a finite weighted Bonferroni estimate with uniform Chinese-remainder counting
errors. Mertens convergence, factorial bounds, and explicit parameter asymptotics
then yield a positive mass--budget gap. All amplified error terms are retained
through the final absorption. The coefficient $1/1000$ is an immediate corollary,
not an optimization target. We describe the ChatGPT--operator workflow, the
formalization's trust boundary, and reusable components, and identify possible
connections to other Erd\H{o}s problems without claiming new results for them.
The supplied final verification run passes all 30 configured gates. The theorem
is a sparsity result: it proves neither finiteness nor the nonexistence of
candidates beyond $24$.
\end{abstract}
\noindent\textbf{Keywords.} Divisor function; consecutive integers; upper-bound sieve;
Bonferroni inequalities; Chinese remainder theorem; Mertens theorem; Lean.

\newpage
\tableofcontents
\newpage

\section{The problem and its arithmetic form}\label{sec:problem}
For a positive integer $m$, let $\tau(m)$ be the number of positive divisors of
$m$, let $\omega(m)$ count its distinct prime divisors, and let $\Omega(m)$ count
its prime factors with multiplicity. All logarithms in this paper are natural.
Erd\H{o}s's 1979 discussions of unconventional problems in number theory include
the question now catalogued as problem \#647, posed with Selfridge
\citep{Erdos1979Magazine,Erdos1979Acta,Bloom647}:
\begin{quote}
Is there an integer $n>24$ such that
$\max_{1\le m<n}\bigl(m+\tau(m)\bigr)\le n+2$?
\end{quote}
The value $n=24$ satisfies the inequality; the question asks for an example
beyond it. Substituting $m=n-k$ makes the same condition a collection of bounds
on a backward window of consecutive integers.

\begin{definition}\label{def:candidate}
A \emph{candidate} is a natural number $n>24$ such that
\begin{equation}\label{eq:candidate}
 \tau(n-k)\le k+2\qquad(1\le k<n).
\end{equation}
Write $\Cand$ for this set, and for $X\in\N$ put
$C(X)=\#(\Cand\cap[1,X])$.
\end{definition}
Here ``candidate'' means an integer satisfying the \emph{entire} condition,
not merely one surviving a finite computation. The first restrictions are
$\tau(n-1)\le3$, $\tau(n-2)\le4$, and $\tau(n-3)\le5$; every further shift is
also constrained. For the proof we may retain only the first $H$ restrictions,
provided $n>H$. This loses information, but gives an upper bound for the actual
candidate count.

\subsection{Position relative to existing work}
Finite exclusion and asymptotic counting address different parts of the problem.
Mian and Siddique give a Lean exclusion certificate for
$24<n\le10^9$, with an audited axiom closure and additional replay checks
reported in their work \citep{MianSiddique2026}. Our argument is not a larger
exclusion computation and does not subsume that finite assertion: its onset is
not numerically specified.

Hughes's consolidation manuscript states both a seven-form bound
$C(X)\ll X/(\log X)^7$ and a growing-window bound
\begin{equation}\label{eq:hughes}
 C(X)\le X\exp\!\bigl(- (\log\log X)^{2-o(1)}\bigr),
\end{equation}
with proofs of those density statements attributed there to companion
manuscripts \citep{Hughes2026}. The saving in the present theorem is
asymptotically larger than the saving displayed in \eqref{eq:hughes}.
Our arithmetic starting point is the logarithmic bound
$\omega(m)\le\log\tau(m)/\log2$, together with an explicit debit for small
primes. This differs from the quadratic-in-$H$ budget obtained by summing the
weaker implication $\Omega(m)+1\le\tau(m)$ over the window.
The comparison is with the stated result in that inspected manuscript; it is
not a claim of an exhaustive priority determination.

Broader work on prime factors of consecutive integers includes the probabilistic
and sieve arguments of Tao--Ter\"av\"ainen and Lau
\citep{TaoTeravainen2026,Lau2026}. We return to specific related problems in
Section~\ref{sec:other}. The classical background is upper-bound sieve theory
and weighted inclusion--exclusion \citep{FriedlanderIwaniec2010,GalambosSimonelli1996}.
The components used here are classical in origin. The contribution described
by this paper is their constant-sensitive assembly for \eqref{eq:candidate},
the resulting bound, and a checked implementation with explicit dependency
boundaries. We do not assert optimality of the exponent or of the numerical
coefficient.

\section{The proved result and its scope}\label{sec:result}
Set
\begin{equation}\label{eq:constants}
 a=\frac{\log2}{1+\log2},\qquad
 \rate=\frac{1499}{10^6},\qquad
 S(X)=\frac{(\log X)^a}{\log\log X}\quad(X\ge3).
\end{equation}
Thus $0<a<1$ and $a\approx0.4094$.

\begin{theorem}[Fixed-coefficient sparsity]\label{thm:main}
For every fixed real $c$ satisfying $0<c<\rate$, there is a natural number
$X_0=X_0(c)\ge3$ such that, for every natural $X\ge X_0$,
\begin{equation}\label{eq:main}
 \boxed{C(X)\le X\exp\!\left(-c\,\frac{(\log X)^a}{\log\log X}\right).}
\end{equation}
In particular, \eqref{eq:main} holds with $c=1/1000$.
\end{theorem}
The main Lean declaration is
\leanname{Erdos647Sieve.Endpoint.endpointBound_of_lt_principal_rate}.
It actually permits any real $c<\rate$; the positive range is the useful
sparsity statement. The explicit corollary is
\leanname{Erdos647Sieve.Endpoint.endpointBound_one_div_thousand};
the original fixed-coefficient specification is discharged by
\leanname{Erdos647Sieve.endpoint} \citep{SieveArtifact}.

The quantifiers are essential: $c$ is selected \emph{before} the onset $X_0$
and the variable $X$. The theorem does not allow a coefficient $c(X)$ tending
to zero, does not provide a numerical $X_0$, and does not give one common onset
for all $c$ approaching $\rate$. Its coefficient range is open at $\rate$.
Section~\ref{sec:absorption} explains why the present majorant needs strict
slack.

\begin{corollary}[Consequences, not an existence decision]\label{cor:density}
The candidate set has natural density zero. For every fixed $M>0$,
$C(X)\le X/(\log X)^M$ for all sufficiently large natural $X$.
Nevertheless, the right-hand side of \eqref{eq:main} tends to infinity.
\end{corollary}
\begin{proof}
Put $Q=\log X$ and $L=\log Q$. Then $S=Q^a/L\to\infty$, and
$S/L=Q^a/L^2\to\infty$. Consequently $e^{-cS}\le e^{-ML}=Q^{-M}$
eventually. On the other hand, $S/Q=Q^{a-1}/\log Q\to0$, so
$\log(Xe^{-cS})=Q-cS\to\infty$.
\end{proof}
These elementary consequences explain the scope of the theorem; they are not
advertised as additional exported declarations in the recorded release.
Even an infinite sufficiently sparse set is compatible with the bound. In
particular, neither the existence nor the nonexistence of one candidate beyond
$24$ follows.

The proof has three levels. A \emph{finite}, translation-uniform moment bound
uses only elementary arithmetic and finite combinatorics. An \emph{analytic}
layer evaluates the available prime mass as the parameters grow. An
\emph{absorption} layer keeps all error terms and converts the resulting
multi-term inequality into \eqref{eq:main}. This separation is reflected in the
Lean packages, rather than being merely expository.

\section{Methodology: AI-assisted construction and formal checking}\label{sec:method}
\subsection{Author, operator, and workflow}
The development was conducted through ChatGPT at \url{https://chatgpt.com}
\citep{OpenAIChatGPT}. The operator reports using the interface label
``GPT-6 Pro'', referred to in this project as Astra. We use the requested author
attribution \emph{GPT 6 Astra}, with \emph{Justin V. Wick} designated as
\emph{operator}. This is an independent project, with no institutional
affiliation. The interface label is workflow metadata supplied by the operator,
not a reproducibly pinned model binary.

The conversational assistant proposed arguments, inspected library interfaces,
produced Lean source and packaging changes, and revised them in response to
compiler output. The operator selected the research direction, ran the project
on a local machine, returned complete diagnostic archives, and managed the
repository and pull-request boundaries. The workflow was iterative rather than
an autonomous, single-call proof search. We make no controlled experimental
claim about model success rates, compute efficiency, or reproducibility of the
same conversation from a prompt.

The mathematical assurance is not the language model's confidence. Proposed
identities and inequalities were expressed as Lean declarations, compiled with
pinned dependencies, and checked against explicit statement and axiom targets.
The ordinary-language proof below is a reconstruction of those ingredients;
it should be read together with the theorem-to-source correspondence in
Appendix~\ref{app:map}. The recorded artifacts, rather than an attempt to replay
the model's internal generation, are the reproducibility object.

\subsection{Reuse before reconstruction}
The formalization uses Lean~4 and Mathlib \citep{Lean4,Mathlib2020,MathlibPin}.
Existing results on factorials, finite Euler products, harmonic sums, floor and
ceiling limits, real powers, and integration were reused where their exact
hypotheses fit. Project-specific adapters handle such issues as real casts of
factorials, integer division versus real division, and equality between a named
function and the expression appearing beneath a derivative.

The mathematical specification was kept separate from implementation repairs.
For example, $B_H$ is signed throughout, repeated weight values are not
identified as indices, and translated intervals may begin at a negative
integer. Those are statement-level choices. Changing a tactic or a module path
does not license changing any of them.

A toolchain migration also exposed a practical limit of indiscriminate reuse:
Lean, Mathlib, and upstream projects must agree on their dependency versions.
The final artifact uses PNT+'s compatible Lean/Mathlib~4.32.2 environment.
The finite package has only Mathlib dependencies. The analytic package isolates
its PNT+ adaptation behind one compatibility module. Thus a consumer of the
finite theorem does not inherit the analytic dependency.

\subsection{What the verification evidence establishes}\label{sec:trust}
The supplied v48 record \citep{SieveArtifact} reports 30 successful configured
gates, 235 distinct declarations with transitive axiom reports, and 220
exact-type or definition examples in 32 audit files. All audited axiom sets are
contained in
\begin{equation}\label{eq:axioms}
 \{\texttt{propext},\ \texttt{Classical.choice},\ \texttt{Quot.sound}\}.
\end{equation}
The successful run has no build or audit warnings, and the recorded source
snapshots are stable. Appendix~\ref{app:evidence} gives the exact artifact and
version identifiers.

These checks have distinct jobs. Compilation constructs and checks proof terms;
expanded-type examples connect the terms to the intended mathematical
statements; axiom inspection detects an admitted or otherwise unauthorized
premise in their transitive dependencies. A theorem of the form
``assuming the moment bound, the conclusion follows'' would not count as a
proof of the moment bound itself. Ordinary runner tests instead check
engineering behavior and do not establish universal mathematical assertions.

The PNT+ repository contains unfinished material outside the slice used here.
The retained 14-module PNT+ import subtree was separately scanned for admitted
source, and the consumed PNT and integrability declarations were axiom-audited.
The unfinished \texttt{ZetaSummary} catalogue is not imported. This is a claim
about the retained dependency, not a declaration that the whole upstream
repository is admission-free \citep{PNTPlus}.

The evidence is an operator-executed build, exact-type, and transitive-axiom
audit of the supplied snapshot. An independently implemented kernel replay and
external mathematical peer review are separate validation levels and are not
claimed for this project. This qualification applies once to the artifact as
a whole; it is not an additional hypothesis in any theorem below.

\section{From divisor restrictions to a signed logarithmic budget}\label{sec:budget}
Fix a natural $H\ge1$ and a real $y>H$. Define
\begin{align}
 \Pset(H,y)&=\{p\text{ prime}:H<p\le y\},\label{eq:prime-set}\\
 F_H(n)&=\prod_{k=1}^H(n-k),\qquad
 T_{H,y}(n)=\#\{p\in\Pset(H,y):p\mid F_H(n)\},\label{eq:hit}\\
 Q_H&=\sum_{k=1}^H\floor{\frac{\log(k+2)}{\log2}},\qquad
 D_H=\sum_{\substack{p\le H\\p\text{ prime}}}\floor{\frac Hp},\qquad
 B_H=Q_H-D_H.\label{eq:budget-def}
\end{align}
We often write $T(n)$ when $H,y$ are fixed. The subtraction defining $B_H$ is
in $\Z$, not truncated subtraction in $\N$.

\begin{lemma}[The occurrence budget]\label{lem:debit}
If $n>H$ and $\tau(n-k)\le k+2$ for $1\le k\le H$, then
\begin{equation}\label{eq:budget}
 T_{H,y}(n)\le B_H.
\end{equation}
More generally, for every integer $n>H$,
\begin{equation}\label{eq:incidence}
 D_H+T_{H,y}(n)\le\sum_{k=1}^H\omega(n-k).
\end{equation}
\end{lemma}
\begin{proof}
Every subset of the distinct prime divisors of a positive integer $m$ gives a
distinct divisor, so $2^{\omega(m)}\le\tau(m)$. Taking logarithms and using the
integrality of $\omega(m)$ gives
\[
 \omega(n-k)\le\floor{\frac{\log(k+2)}{\log2}}.
\]
Thus the total number of prime--shift incidences is at most $Q_H$.

For a prime $p\le H$, every $H$ consecutive integers contain at least
$\floor{H/p}$ multiples of $p$. These mandatory incidences contribute at least
$D_H$. Each selected prime counted by $T(n)$ divides at least one factor of
$F_H(n)$ and contributes a further incidence. The small-prime and selected-prime
sets are disjoint. Counting all these incidences proves \eqref{eq:incidence};
combining it with the preceding upper bound proves \eqref{eq:budget}.
\end{proof}
Since $T(n)\ge0$, a negative value of $B_H$ rules out every such prefix survivor
above $H$. In particular,
\begin{equation}\label{eq:negative-budget}
 B_H<0\quad\Longrightarrow\quad C(X)\le H.
\end{equation}
This simple branch must remain available. Proving eventual nonnegativity of
the corrected budget is unnecessary for the endpoint.

The key improvement is the size of the available budget. Without the debit,
$Q_H$ is of order $H\log H$. Small primes consume order $H\log\log H$
occurrences regardless of where the window lies. Subtracting that forced cost
is precisely what permits constant-level cancellation at the critical exponent
chosen later. The proof is counting prime occurrences, not assuming independent
factorization of neighboring integers.

\section{A finite weighted moment estimate}\label{sec:moment}
The following theorem is the finite core. It is independent of the prime number
theorem, Mertens convergence, and the later parameter choices.

\begin{theorem}[Uniform translated finite moment]\label{thm:moment}
Let $A\in\Z$, let $X,H\ge1$ be natural numbers, let $y>H$, let $0<z<1$, and
let $J\ge0$ be an even natural number. Put
\begin{equation}\label{eq:mass}
 t=1-z,\qquad \lambda=H\sum_{p\in\Pset(H,y)}\frac1p,\qquad \mu=t\lambda.
\end{equation}
Then
\begin{equation}\label{eq:moment}
 \sum_{A<n\le A+X}z^{T_{H,y}(n)}
 \le X\left(e^{-\mu}+\frac{\mu^{J+1}}{(J+1)!}\right)+(Hy)^J.
\end{equation}
\end{theorem}
The summation in \eqref{eq:moment} is over integers. In particular, $A$ may be
negative and the interval may contain a zero of $F_H$. No positive-window
assumption is used here. The bound includes the empty selected-prime set,
$J=0$, and truncation orders exceeding the number of selected primes.

\subsection{Weighted Bonferroni algebra}
For a finite index set $P$ and real weights $(w_p)_{p\in P}$, write
\[
 e_j(w)=\sum_{\substack{S\subseteq P\\|S|=j}}\prod_{p\in S}w_p,
 \qquad e_0(w)=1,
 \qquad P_J(w)=\sum_{j=0}^J(-1)^j e_j(w).
\]
Empty products equal one; coefficients of degree greater than $|P|$ are zero.
The indices are the elements of $P$. Equal values of $w_p$ are not deduplicated.

\begin{lemma}[Three finite inequalities]\label{lem:bonferroni}
For $0\le w_p\le1$ and even $J$,
\begin{equation}\label{eq:bonferroni}
 \prod_{p\in P}(1-w_p)\le P_J(w)
 \le\prod_{p\in P}(1-w_p)+e_{J+1}(w).
\end{equation}
For nonnegative weights and every natural $j$,
\begin{equation}\label{eq:coefficient}
 j!\,e_j(w)\le\left(\sum_{p\in P}w_p\right)^j.
\end{equation}
Finally, if $0\le w_p\le1$, then
\begin{equation}\label{eq:product-exp}
 \prod_{p\in P}(1-w_p)\le\exp\!\left(-\sum_{p\in P}w_p\right).
\end{equation}
\end{lemma}
\begin{proof}
The weighted Bonferroni inequality \eqref{eq:bonferroni} follows from finite
inclusion--exclusion; Appendix~\ref{app:bonferroni} gives the recurrence used in
the formalization. This is the usual deterministic weighted inequality, not a
statement of independence for the arithmetic events \citep{GalambosSimonelli1996}.
For \eqref{eq:coefficient}, expand the power on the right as a sum over ordered
$j$-tuples. The tuples with distinct indices contribute exactly $j!e_j(w)$,
and all other terms are nonnegative. The case $j=0$ gives equality.
For \eqref{eq:product-exp}, multiply $1-u\le e^{-u}$ over the indices, using
nonnegativity of all factors.
\end{proof}
For the arithmetic weights and their model values, set
\begin{equation}\label{eq:weights}
 w_p(n)=t\ind_{p\mid F_H(n)},\qquad b_p=\frac{tH}{p}
 \quad(p\in\Pset(H,y)).
\end{equation}
Both sets of weights lie in $[0,1]$, and
\begin{equation}\label{eq:weight-products}
 \prod_p(1-w_p(n))=z^{T(n)},\qquad \sum_p b_p=\mu.
\end{equation}
Applying the first half of \eqref{eq:bonferroni} pointwise gives
\begin{equation}\label{eq:pointwise-summed}
 \sum_{A<n\le A+X}z^{T(n)}
 \le\sum_{j=0}^J(-1)^j\sum_{A<n\le A+X}e_j(w(n)).
\end{equation}
Applying the other half to $b$, together with
\eqref{eq:coefficient}--\eqref{eq:product-exp}, gives
\begin{equation}\label{eq:mean-truncation}
 P_J(b)\le e^{-\mu}+\frac{\mu^{J+1}}{(J+1)!}.
\end{equation}
It remains to compare the actual coefficient sums in
\eqref{eq:pointwise-summed} with $XP_J(b)$. The name ``model value'' does not
supply this comparison; the next two steps prove it arithmetically.

\subsection{Root counting and arbitrary translated intervals}
For a subset $S\subseteq\Pset(H,y)$, define
\[
 d(S)=\prod_{p\in S}p,\qquad
 N_S=\#\{A<n\le A+X:d(S)\mid F_H(n)\}.
\]
The integer $d(S)$ is positive and squarefree, including $d(\varnothing)=1$.
For distinct primes, simultaneous divisibility is equivalent to divisibility by
the product, also when $F_H(n)=0$.

\begin{lemma}[Uniform subset-count discrepancy]\label{lem:crt}
For every such $S$,
\begin{equation}\label{eq:subset-error}
 \left|N_S-\frac{XH^{|S|}}{d(S)}\right|\le H^{|S|}.
\end{equation}
\end{lemma}
\begin{proof}
For a prime $p>H$, the congruence $F_H(n)\equiv0\pmod p$ holds exactly in the
$H$ distinct residue classes $1,\ldots,H$. The Chinese remainder theorem
therefore gives precisely $H^{|S|}$ root classes modulo $d(S)$.

For every positive integer modulus $d$ and integer representative $r$, the
number of points of that residue class in the translated interval is exactly
\begin{equation}\label{eq:floor-count}
 \#\{A<n\le A+X:d\mid n-r\}
 =\floor{\frac{A+X-r}{d}}-\floor{\frac{A-r}{d}}.
\end{equation}
Indeed the points have the form $r+jd$, and their allowed integer indices form
the half-open interval between the two displayed quotients. Subtracting $X/d$
from \eqref{eq:floor-count} leaves the difference of two fractional parts,
whose absolute value is at most one. Summing the error over the distinct root
classes proves \eqref{eq:subset-error}. For $S=\varnothing$, the modulus is
one, there is one residue class, and $N_S=X$; the same inequality applies.
\end{proof}
The floor formula is valid for negative $A$ and $r$, because its divisions are
real quotients followed by floors, corresponding to Euclidean integer division
in Lean. No truncation-toward-zero convention is being used.

\subsection{Summing every retained subset once}
The subset interpretation of the elementary coefficients gives
\begin{align}
 \sum_{A<n\le A+X} e_j(w(n))
 &=t^j\sum_{\substack{S\subseteq\Pset(H,y)\\|S|=j}}N_S,\label{eq:actual-coeff}\\
 e_j(b)&=t^j\sum_{\substack{S\subseteq\Pset(H,y)\\|S|=j}}
                  \frac{H^j}{d(S)}.\label{eq:model-coeff}
\end{align}
Consequently their complete signed discrepancy is
\begin{align}
 \mathcal E_J
 &:=\sum_{j=0}^J(-1)^j
       \left(\sum_{A<n\le A+X}e_j(w(n))-Xe_j(b)\right)\notag\\
 &=\sum_{\substack{S\subseteq\Pset(H,y)\\|S|\le J}}
          (-t)^{|S|}\left(N_S-\frac{XH^{|S|}}{d(S)}\right).\label{eq:discrepancy}
\end{align}
We take absolute values only after writing this signed identity. Multiplying
unrelated one-sided estimates by $(-1)^j$ would not be valid.

Unique factorization makes $S\mapsto d(S)$ injective. Since $y>H\ge1$, every
retained subset satisfies
\[
 1\le d(S)\le y^{|S|}\le y^J.
\]
Hence the entire retained family has at most $\floor{y^J}$ members. From
\eqref{eq:subset-error}, $0<t<1$, and $H\ge1$,
\begin{equation}\label{eq:aggregate}
 |\mathcal E_J|
 \le\sum_{\substack{S\subseteq\Pset(H,y)\\|S|\le J}}t^{|S|}H^{|S|}
 \le H^J\floor{y^J}\le(Hy)^J.
\end{equation}
This counts each subset once across all degrees. It is the reason the error has
no additional factor $J+1$. For $J=0$ the only subset is empty, and the bound
remains valid.

\begin{proof}[Completion of Theorem~\ref{thm:moment}]
Use \eqref{eq:pointwise-summed}, write its right-hand side as $XP_J(b)+\mathcal E_J$,
and apply \eqref{eq:mean-truncation} and \eqref{eq:aggregate}.
\end{proof}

\section{Transferring the moment bound to actual candidates}\label{sec:transfer}
Let $\eta=-\log z>0$. For a candidate $n>H$, Lemma~\ref{lem:debit} gives
$T(n)\le B_H$. If $B_H\ge0$, monotonicity of $z^u$ for $0<z<1$ implies
\[
 1\le e^{\eta B_H}z^{T(n)}.
\]
Counting at most $H$ candidate integers in the initial window, and applying
Theorem~\ref{thm:moment} with $A=0$, gives the following exact finite transfer.

\begin{theorem}[Finite candidate-counting inequality]\label{thm:finite-count}
Under the numerical hypotheses of Theorem~\ref{thm:moment}, if $B_H\ge0$ then
\begin{equation}\label{eq:finite-count}
 C(X)\le H+e^{\eta B_H}
 \left\{X\left(e^{-\mu}+\frac{\mu^{J+1}}{(J+1)!}\right)+(Hy)^J\right\}.
\end{equation}
If $B_H<0$, the separate bound $C(X)\le H$ holds.
\end{theorem}
The proof uses a necessary condition for every actual candidate, so the count
is not limited to an algorithmic survivor set. Also, the multiplier
$e^{\eta B_H}$ applies to the \emph{entire} expression inside braces.
Discarding it from the factorial tail or the arithmetic remainder would change
the theorem. These amplified contributions will be treated separately rather
than hidden in an error term.

The two finite main declarations are
\leanname{Erdos647Sieve.finiteMoment} and
\leanname{Erdos647Sieve.finiteCounting}. They have no unproved moment, CRT,
or distribution assertion as a premise. Their displayed numerical hypotheses,
of course, remain part of their types \citep{SieveArtifact}.

\section{Elementary estimates with the linear term retained}\label{sec:elementary}
The budget comparison requires more than its leading order. This section keeps
the constant contribution coming from the linear term in $\log(H!)$.

\subsection{The mandatory small-prime debit}
\begin{lemma}[An elementary reciprocal-prime lower bound]\label{lem:prime-lower}
For every natural $H\ge2$,
\begin{equation}\label{eq:prime-lower}
 \sum_{p\le H}\frac1p\ge\log\log H-1.
\end{equation}
Here and below an index $p$ denotes a prime.
\end{lemma}
\begin{proof}
Expanding the finite-prime-base Euler product gives
\[
 E_H:=\prod_{p\le H}(1-1/p)^{-1}
   =\sum_{\substack{m\ge1\\\text{all prime factors of }m\le H}}\frac1m
   \ge\sum_{m=1}^H\frac1m\ge\log H.
\]
The sum over smooth integers converges because only finitely many prime bases
are involved and every prime-power geometric series converges. This does not
assert convergence of the full harmonic series. Positivity permits the
restriction to $1\le m\le H$.
For $0\le u<1$, integration of $v/(1-v)$ gives
\[
 -\log(1-u)-u\le\frac{u^2}{1-u}.
\]
At $u=1/p$, the right-hand side is $1/[p(p-1)]$. The error over primes is bounded
by the telescoping sum
\[
 \sum_{p\le H}\frac1{p(p-1)}
 \le\sum_{m=2}^H\left(\frac1{m-1}-\frac1m\right)
 =1-\frac1H\le1.
\]
Taking logarithms of the Euler-product comparison and summing the factor bounds
yields $\log\log H\le\log E_H\le\sum_{p\le H}1/p+1$.
\end{proof}

\begin{corollary}[Fixed-constant debit]\label{cor:debit-lower}
For every natural $H\ge2$,
\begin{equation}\label{eq:debit-lower}
 D_H\ge H(\log\log H-2).
\end{equation}
\end{corollary}
\begin{proof}
Use $\floor{H/p}\ge H/p-1$, sum over the primes at most $H$, and bound their
number by $H$. Lemma~\ref{lem:prime-lower} then proves the claim.
\end{proof}
The lower bound can be negative for small $H$; this is harmless. In the
formalization the quotients in $D_H$ are natural-number quotients cast to
$\R$ only after division. The comparison with real division is a proved
rounding lemma.

\subsection{Factorials and the uncorrected budget}
The proof uses existing Mathlib Stirling machinery, not a new formalization of
Stirling's formula \citep{MathlibPin}. Its convenient consequences are
\begin{equation}\label{eq:stirling}
 n\log n-n\le\log(n!)
 \le n\log n-n+\frac12\log n+1\qquad(n\ge1),
\end{equation}
and
\begin{equation}\label{eq:factorial-lower}
 (n/e)^n\le n!\qquad(n\ge1).
\end{equation}
For instance, the upper estimate follows from monotonicity of the Stirling
sequence $n!/[\sqrt{2n}(n/e)^n]$ and its value at $n=1$; its lower bound
$\sqrt\pi$ implies the lower estimate in \eqref{eq:stirling} after dropping
nonnegative logarithmic terms. Only these consequences are needed here.

Put
\[
 L_{\mathrm{fac}}(H)=\log\left(\frac{(H+2)!}{2}\right).
\]
The division is real division after casting the factorial. The exact identity
\begin{equation}\label{eq:factorial-shift}
 L_{\mathrm{fac}}(H)=\log(H!)+\log(H+1)+\log(H+2)-\log2
\end{equation}
and \eqref{eq:stirling} give
\begin{equation}\label{eq:factorial-limit}
 \frac{L_{\mathrm{fac}}(H)-(H\log H-H)}{H}\longrightarrow0
 \qquad(H\in\N,\ H\to\infty).
\end{equation}
On the other hand, bounding each floor by its argument and summing logarithms
of the positive factors $3,\ldots,H+2$ gives
\begin{equation}\label{eq:log-budget}
 Q_H\le\frac1{\log2}\sum_{k=1}^H\log(k+2)
      =\frac{L_{\mathrm{fac}}(H)}{\log2}.
\end{equation}

\begin{proposition}[One-sided corrected-budget estimate]\label{prop:budget-upper}
For every $\varepsilon>0$, eventually over natural $H$,
\begin{equation}\label{eq:budget-upper}
 \frac{B_H}{H}\le
 \frac{\log H}{\log2}-\log\log H+2-\frac1{\log2}+\varepsilon.
\end{equation}
\end{proposition}
\begin{proof}
Subtract \eqref{eq:debit-lower} from \eqref{eq:log-budget}, divide by $H>0$,
and use \eqref{eq:factorial-limit}.
\end{proof}
The term $-1/\log2$ is the normalized contribution of $-H$ in the factorial
estimate. Replacing that estimate by $\log(H!)\le H\log H$ would lose this
constant. We do not assert an asymptotic equality for $B_H/H$: integer floors
can contribute a linear total error, and the proof only needs the indicated
one-sided bound.

\section{The analytic input: convergent Mertens error}\label{sec:analytic}
Write
\[
 \vartheta(x)=\sum_{p\le x}\log p,\quad
 E(x)=\vartheta(x)-x,\quad
 R(x)=\sum_{p\le x}\frac1p-\log\log x\qquad(x\ge2).
\]
The classical reciprocal-prime analysis is closely related to the
Rosser--Schoenfeld formulas \citep{RosserSchoenfeld1962}. The formal input used
here is supplied by a restricted, audited portion of PNT+
\citep{PNTPlus,SieveArtifact}:
\begin{equation}\label{eq:pnt}
 \exists C\ge0\quad\forall x\ge2,\qquad
 |E(x)|\le\frac{Cx}{(\log x)^2}.
\end{equation}
The source proves this through its \texttt{MediumPNT} theorem. No numerical
zero verification or zero-free-region catalogue from \texttt{ZetaSummary} is
used by the retained route.

\begin{proposition}[Mertens convergence in the required form]\label{prop:mertens}
There is a real number $M$ such that $R(x)\to M$ as $x\to\infty$.
Consequently, if $u,v$ are two real-valued cutoffs tending to infinity, then
\begin{equation}\label{eq:mertens-cancel}
 \sum_{p\le v}\frac1p-\sum_{p\le u}\frac1p
 -\bigl(\log\log v-\log\log u\bigr)\longrightarrow0.
\end{equation}
\end{proposition}
\begin{proof}
For $f(x)=1/(x\log x)$ on $(1,\infty)$,
\[
 f'(x)=-\frac{1+\log x}{x^2(\log x)^2}.
\]
Specializing Abel summation to the prime weights $\log p$ gives the exact
identity
\begin{equation}\label{eq:abel}
 \sum_{p\le x}\frac1p
 =\frac{\vartheta(x)}{x\log x}
   +\int_2^x\frac{\vartheta(u)(1+\log u)}{u^2(\log u)^2}\,du.
\end{equation}
For completeness, the lower endpoint is consistent at $x=2$: the boundary
term is $1/2$ and the integral is zero. There is no omitted endpoint constant.
Writing $\vartheta(u)=u+E(u)$ and using the antiderivative
\[
 G(u)=\frac1{\log u}-\log\log u,\qquad G'(u)=u f'(u),
\]
rearranges \eqref{eq:abel} to
\begin{equation}\label{eq:mertens-error}
 R(x)=G(2)+\frac{E(x)}{x\log x}-\int_2^x E(u)f'(u)\,du.
\end{equation}
By \eqref{eq:pnt}, the boundary term tends to zero and
\[
 |E(u)f'(u)|\le C\left(\frac1{u(\log u)^3}
                         +\frac1{u(\log u)^4}\right).
\]
The latter majorant is integrable on $[2,\infty)$. Thus the right-hand side of
\eqref{eq:mertens-error} converges, proving existence of $M$.
Composing this convergence with $u$ and $v$ and subtracting proves
\eqref{eq:mertens-cancel}.
\end{proof}
The actual value of $M$ is irrelevant: it cancels. A merely bounded error
$R(x)=O(1)$ would leave an uncontrolled constant in \eqref{eq:mertens-cancel}
and would not suffice for the later positive gap.

For real $y\ge H$, the exact selected-prime identity is
\begin{equation}\label{eq:prime-difference}
 \sum_{p\in\Pset(H,y)}\frac1p
 =\sum_{p\le y}\frac1p-\sum_{p\le H}\frac1p.
\end{equation}
The formalization proves this finite-set partition and then the corresponding
limit for the actual selected-prime definition. The PNT+ constants and
interfaces are hidden behind the analytic compatibility boundary; downstream
parameter proofs consume \eqref{eq:mertens-cancel}, not upstream internals.

\section{Parameter asymptotics and the positive mass--budget gap}\label{sec:parameters}
For the remainder of the proof $X$ tends to infinity through natural numbers.
Use the fixed choices
\begin{equation}\label{eq:parameters}
 \begin{gathered}
 Q=\log X,\qquad L=\log Q,\qquad
 H=\floor{Q^a},\qquad
 J=2\ceil{\frac H{100}},\\
 y=X^{1/(4J)},\qquad t=\frac1{1000L},\qquad z=1-t.
 \end{gathered}
\end{equation}
These are the existing formal parameter definitions. In particular, the natural
expression for $J$ is $2((H+99)/100)$, with natural division. Its equality to the
ceiling expression in \eqref{eq:parameters} is proved separately.

\subsection{Why this exponent is selected}
If a provisional window were $H\asymp Q^b$ and $J\asymp H$, then
$\log y\asymp Q^{1-b}$. The leading prime mass per shift would have coefficient
$1-b$ in front of $L$, while the logarithmic budget would have coefficient
$b/\log2$. The equality
\begin{equation}\label{eq:critical}
 \frac a{\log2}=1-a
\end{equation}
selects $a=\log2/(1+\log2)$. At this boundary the leading logarithms cancel,
so the secondary terms and their fixed constants matter. This identifies the
critical balance of the chosen argument, not an optimality theorem among all
possible methods.

\begin{lemma}[Actual window, truncation, and cutoff limits]\label{lem:parameters}
For the parameters in \eqref{eq:parameters},
\begin{align}
 &H\to\infty,\qquad H/Q^a\to1,\qquad\log H-aL\to0,\label{eq:window-limits}\\
 &H/50\le J<H/50+2,\qquad J/H\to1/50,\label{eq:truncation-limits}\\
 &\frac{\log y}{Q^{1-a}}\to\frac{25}{2},\qquad
   \log\log y-(1-a)L\to\log(25/2).\label{eq:cutoff-limits}
\end{align}
Moreover $y\to\infty$, $H<y$ eventually, $J$ is even, and all numerical
hypotheses of Theorem~\ref{thm:moment} hold eventually.
\end{lemma}
\begin{proof}
The inequalities $Q^a-1<H\le Q^a$ prove the first two assertions in
\eqref{eq:window-limits}. For sufficiently large $X$ the quantities are
positive, and
$\log H-aL=\log(H/Q^a)\to0$.
The exact ceiling identity for $J$ proves \eqref{eq:truncation-limits}.

For positive $X$ and $J$, the real-power definition gives
$\log y=Q/(4J)$. Hence
\[
 \frac{\log y}{Q^{1-a}}=\frac{Q^a}{4J}\longrightarrow\frac{25}{2}.
\]
Taking logarithms proves the second assertion in \eqref{eq:cutoff-limits}.
The power $Q^{1-a}$ tends to infinity and dominates $L$. Therefore
$\log y\to\infty$ and $\log H/\log y\to0$, so $y\to\infty$ and $H<y$
eventually. Finally, $L\to\infty$ implies $0<t\le1/2$ eventually, and thus
$0<z<1$. Positivity of $H,J$ and evenness of $J$ have already been established.
\end{proof}
These arguments use familiar limits, but their formal application requires the
floor, ceiling, cast, and nonzero-denominator obligations. They are supplied
for these specific functions of $X$, rather than left as cutoff hypotheses.

\begin{proposition}[Prime-mass expansion]\label{prop:mass}
Let $\lambda=\lambda(H,y)$ for \eqref{eq:parameters}. Then
\begin{equation}\label{eq:mass-expansion}
 \frac{\lambda}{H}
 = (1-a)L-\log\log H+\log(25/2)+o(1).
\end{equation}
\end{proposition}
\begin{proof}
The two cutoffs $H,y$ tend to infinity and satisfy $H\le y$ eventually.
Apply \eqref{eq:mertens-cancel} and \eqref{eq:prime-difference} to obtain
$\lambda/H=\log\log y-\log\log H+o(1)$, and then use
\eqref{eq:cutoff-limits}.
\end{proof}

\begin{proposition}[Fixed positive gap and size bounds]\label{prop:gap}
Define
\begin{equation}\label{eq:K}
 K=\log(25/2)+\frac1{\log2}-2.
\end{equation}
For every fixed real $\delta<K$, eventually
$\lambda-B_H\ge\delta H$. In particular,
\begin{equation}\label{eq:gap-and-bounds}
 \lambda-B_H\ge\frac32H,\qquad
 0\le\lambda\le HL,\qquad B_H\le HL
 \quad\text{eventually}.
\end{equation}
Also $\lambda/(HL)\to1-a$.
\end{proposition}
\begin{proof}
Let
\[
 F(H)=\frac{\log H}{\log2}-\log\log H+2-\frac1{\log2}.
\]
From \eqref{eq:mass-expansion}, \eqref{eq:window-limits}, and
\eqref{eq:critical},
\begin{equation}\label{eq:gap-limit}
 \frac\lambda H-F(H)\longrightarrow K.
\end{equation}
By Proposition~\ref{prop:budget-upper}, for every $\varepsilon>0$ one has
$B_H/H\le F(H)+\varepsilon$ eventually, also after composing with $H(X)$.
Given $\delta<K$, choose $\varepsilon=(K-\delta)/3$.
Eventually \eqref{eq:gap-limit} is at least $K-\varepsilon$, so
$(\lambda-B_H)/H\ge K-2\varepsilon>\delta$.
The certified elementary logarithm estimate $K>3/2$ is recorded in
Appendix~\ref{app:constants}.

For the size bound, $\log H/L\to a$, and
\[
 \frac{\log\log H}{L}
 =\frac{\log\log H}{\log H}\,\frac{\log H}{L}\longrightarrow0.
\]
Dividing \eqref{eq:mass-expansion} by $L$ gives
$\lambda/(HL)\to1-a<1$. The mass is nonnegative by definition; hence
$0\le\lambda\le HL$ eventually. The positive gap gives $B_H\le\lambda$,
which yields the remaining upper bound.
\end{proof}
There is no assertion that $(\lambda-B_H)/H$ converges to $K$: only
\eqref{eq:gap-limit} has been shown to converge, and the budget estimate is
one-sided. Nor is any sign hypothesis on $B_H$ part of
Proposition~\ref{prop:gap}.

\section{Every amplified contribution}\label{sec:errors}
We now bound the four terms in Theorem~\ref{thm:finite-count}, using
\eqref{eq:parameters}. Write
\[
 \eta=-\log(1-t),\qquad \mu=t\lambda,
 \qquad A_X=\eta B_H.
\]
The symbol $A_X$ in this section is an amplification exponent, not the interval
translation used earlier. For clarity, all inequalities involving the positive
multiplier are first proved in the case $B_H\ge0$.

\subsection{Principal saving and the source of 1499}
For $0\le t\le1/2$, elementary calculus gives
\begin{equation}\label{eq:log-weight}
 0\le-\log(1-t)\le t+t^2\le2t.
\end{equation}
For example the derivative of $t+t^2+\log(1-t)$ is
$t(1-2t)/(1-t)\ge0$ on this interval, and the function vanishes at zero.
Multiplying \eqref{eq:log-weight} by the nonnegative budget and using
\eqref{eq:gap-and-bounds} gives
\begin{equation}\label{eq:amplification}
 A_X\le2tHL=H/500,
 \qquad 0\le\mu\le tHL=H/1000.
\end{equation}
More importantly,
\begin{align}
 A_X-\mu
 &\le(t+t^2)B_H-t\lambda\notag\\
 &=-t(\lambda-B_H)+t^2 B_H\notag\\
 &\le-\frac32tH+t^2HL
 =-\frac{1499}{10^6}\frac HL.\label{eq:principal}
\end{align}
Thus the principal contribution is at most $Xe^{-\rate H/L}$.
The integer $1499$ is simply $1500-1$: the first part is the gap saving and the
second pays for the quadratic logarithm correction. It is not an invariant of
problem \#647 and is not claimed to be optimal.

\subsection{The factorial tail with a growing numerator}
Put $q=J+1$. By \eqref{eq:truncation-limits}, $q\ge H/50$ and $q\ge1$
eventually. Equation~\eqref{eq:amplification} gives $\mu/q\le1/20$.
Consequently \eqref{eq:factorial-lower} yields
\begin{equation}\label{eq:factorial-tail}
 \frac{\mu^q}{q!}\le\left(\frac{e\mu}{q}\right)^q
 \le(e/20)^q.
\end{equation}
Since $\log20-1>0$ and $q\ge H/50$,
\begin{align}
 e^{A_X}\frac{\mu^q}{q!}
 &\le\exp\!\left(\frac H{500}-(\log20-1)q\right)\notag\\
 &\le\exp\!\left[-H\left(\frac{\log20-1}{50}-\frac1{500}\right)\right]
 \le e^{-H/30}.\label{eq:amplified-tail}
\end{align}
The last rational comparison is included in Appendix~\ref{app:constants}.
This is a bound on the actual growing numerator $\mu(X)$; the elementary limit
$u^q/q!\to0$ for a \emph{fixed} $u$ would not justify this step.

\subsection{Arithmetic remainder and the exceptional window}
For the actual positive parameters, the power laws give the exact identity
\begin{equation}\label{eq:y-power}
 y^J=X^{1/4}.
\end{equation}
Keeping the full multiplier on the discrepancy term,
\begin{equation}\label{eq:arithmetic-log}
 e^{A_X}(Hy)^J
 =\exp\!\left(\frac Q4+J\log H+A_X\right).
\end{equation}
Since $J\asymp H\asymp Q^a$ and $a<1$,
\[
 \frac{J\log H+H/500}{Q}\longrightarrow0.
\]
It follows from \eqref{eq:amplification}--\eqref{eq:arithmetic-log} that
\begin{equation}\label{eq:arithmetic-bound}
 e^{A_X}(Hy)^J\le e^{Q/2}=\sqrt X\qquad\text{eventually}.
\end{equation}
Independently $H\asymp Q^a\le\sqrt X$ eventually, so the exceptional initial
window also contributes at most $\sqrt X$.

\begin{theorem}[All-sign intermediate bound]\label{thm:majorant}
For all sufficiently large natural $X$,
\begin{equation}\label{eq:majorant}
 \boxed{C(X)\le Xe^{-\rate H/L}+Xe^{-H/30}+2\sqrt X.}
\end{equation}
There is no assumption on the sign of $B_H$ in this statement.
\end{theorem}
\begin{proof}
For $B_H\ge0$, apply Theorem~\ref{thm:finite-count} and the three estimates
\eqref{eq:principal}, \eqref{eq:amplified-tail}, and
\eqref{eq:arithmetic-bound}, along with the exceptional-window bound.
For $B_H<0$, \eqref{eq:negative-budget} gives $C(X)\le H\le\sqrt X$, which is
smaller than the displayed nonnegative right-hand side. The finitely many
eventual conditions used in these comparisons can be imposed simultaneously.
\end{proof}
The formal declaration is
\leanname{Erdos647Sieve.Endpoint.eventually_candidateCount_le_amplified_errors}.
It is the last counting input to the final assembly. Every error source from
the finite theorem has a named bound before that assembly is attempted.

\section{Final absorption and the strict coefficient range}\label{sec:absorption}
Let $U=Q^a$ and $S=U/L$. The parameter limits give
\begin{equation}\label{eq:absorption-scales}
 S\to\infty,\quad H/U\to1,\quad S/H\to0,\quad S/Q\to0.
\end{equation}
For example $S/H=(U/H)/L\to0$, and $S/Q=Q^{a-1}/L\to0$.
All denominators are positive eventually.

\begin{proof}[Proof of Theorem~\ref{thm:main}]
Fix $0<c<\rate$. Divide the right-hand side of \eqref{eq:majorant} by the
positive quantity $Xe^{-cS}$. The exact resulting ratio is
\begin{equation}\label{eq:normalized-majorant}
 \exp\!\left[\left(c-\rate\frac HU\right)S\right]
 +\exp\!\left[\left(c\frac SH-\frac1{30}\right)H\right]
 +2\exp\!\left[\left(c\frac SQ-\frac12\right)Q\right].
\end{equation}
By \eqref{eq:absorption-scales}, the three coefficients in parentheses tend
respectively to $c-\rate<0$, $-1/30$, and $-1/2$. Each is eventually bounded
above by a fixed negative number. Their multiplying scales $S,H,Q$ tend to
infinity. Therefore every term of \eqref{eq:normalized-majorant} tends to
zero, and the sum is eventually at most one. Combining this with
Theorem~\ref{thm:majorant} proves \eqref{eq:main}. Finally
$0<1/1000<1499/10^6$ is rational arithmetic, so the stated specialization
requires no additional analytic estimate.
\end{proof}
This proof absorbs the rounding of $H$, the factor two, and both secondary
terms without an unspecified multiplicative constant in the conclusion.
Its limit statements assert the existence of an onset; they do not compute
that onset.

\subsection{Why the boundary coefficient is not asserted}
At $c=\rate$, the first ratio in \eqref{eq:normalized-majorant} becomes
\begin{equation}\label{eq:boundary-ratio}
 \exp\!\left(\rate\frac{U-\floor U}{L}\right)\ge1.
\end{equation}
The other two terms are positive. Thus the available majorant cannot itself
be bounded by $Xe^{-\rate S}$ using this argument. This does not prove that the
actual candidate count violates such a bound; it shows why the established
coefficient range is strict. Taking $c$ to the boundary would also require
control of the dependence of $X_0(c)$ on $c$.

The right presentation is therefore the complete range $0<c<\rate$, followed
by a simple explicit corollary. The parameter $t=1/(1000L)$ and the final
coefficient $c$ have different roles. The former tunes the moment; the latter
is the saving that survives all losses. Here the original $c=1/1000$ target
happens to follow for free, but the proof never needs to optimize that
arbitrary value.

\section{Reusable mathematical and formal machinery}\label{sec:reuse}
The proof yields more reusable material than the final statement alone. This
section separates existing library components from natural extensions that
have not been exported as Lean theorems in the audited snapshot.

\subsection{Components already implemented}
The finite package provides weighted Bonferroni inequalities for finite lists,
coefficient bounds, and a bridge from list coefficients to subsets of a finite
index set. These do not depend on the divisor problem. The distinction between
indices and weight values allows repeated weights without accidental loss of
multiplicity. They are suitable building blocks for deterministic lower-tail
estimates in other finite sieves.

The progression-counting result \eqref{eq:floor-count} is uniform in the
translation and residue representative. The prime-subset product lemmas
provide injective encoding by positive squarefree products. The combination
of these with CRT supplies exact local densities and an explicit cumulative
error. Their useful abstraction is not ``primes behave independently'', but
``simultaneous congruences have exact finite root counts, and interval sampling
has a controlled error''.

The seven elementary modules promoted to the Mathlib-only package expose
factorial bounds and limits, the finite Euler-product comparison, the
reciprocal-prime lower bound, the mandatory debit, and the signed-budget
estimate. Their public facade is \smallcode{Erdos647Sieve.Elementary}.
The finite theorem is accessible through \smallcode{Erdos647Sieve.Finite}.
Neither facade requires PNT+ \citep{SieveArtifact}.

The analytic package supplies Mertens convergence and its two-cutoff
cancellation in the precise form consumed by the application. Its public
facades are \smallcode{Erdos647Research.Analytic} and
\smallcode{Erdos647Research.Endpoint}. The latter exposes the completed
endpoint; ``Research'' in the module name indicates its separate dependency
boundary, not an unfinished proof. Only the compatibility layer knows the
upstream PNT+ interfaces.

\subsection{A general residue-set template}\label{sec:template}
The following is a mathematical extension of the finite argument, stated to
make the transfer mechanism explicit. It is \emph{not} claimed as a separately
exported generic theorem in the v48 Lean package.
Let $P$ be any finite set of distinct primes. At each $p\in P$ choose a residue
set $R_p\subseteq\Z/p\Z$ of size $\rho_p$. Define
\[
 T_R(n)=\sum_{p\in P}\ind_{n\bmod p\in R_p},\qquad
 \rho(S)=\prod_{p\in S}\rho_p,\qquad
 d(S)=\prod_{p\in S}p.
\]
For $0<z<1$, $t=1-z$, even $J\ge0$, and an interval of $X$ consecutive integers,
the same proof gives
\begin{equation}\label{eq:template}
 \sum_{A<n\le A+X}z^{T_R(n)}
 \le X\left(e^{-\mu_R}+\frac{\mu_R^{J+1}}{(J+1)!}\right)
  +\sum_{\substack{S\subseteq P\\|S|\le J}}t^{|S|}\rho(S),
 \qquad \mu_R=t\sum_{p\in P}\frac{\rho_p}{p}.
\end{equation}
Indeed CRT gives $\rho(S)$ simultaneous classes modulo $d(S)$, the counting
error is at most $\rho(S)$, and the model weights $t\rho_p/p$ lie in $[0,1]$.
Equations~\eqref{eq:bonferroni}--\eqref{eq:discrepancy} then give
\eqref{eq:template}. This proof on paper specifies the additional generic
interface one could formalize; it does not silently generalize existing theorem
types. Bounding the last sum still requires information on the particular
residue sets and cutoff products.

\subsection{Other budget profiles and practical limits}
For restrictions $\tau(n-k)\le f(k)$ with $f(k)\ge1$, the same elementary
argument suggests the budget
\[
 \sum_{k=1}^H\floor{\frac{\log f(k)}{\log2}}-D_H.
\]
For a nonconsecutive family of shifts, both the mandatory small-prime debit and
the local root count must be recalculated. If several shifts coincide modulo
$p$, the root density is no longer $H/p$. If prime powers are included as sieve
moduli, prime-subset product injectivity alone is no longer the correct
bookkeeping device. These are genuine extension obligations.

The final ratio-absorption pattern is another broadly useful component: prove
a normalized error tends to zero, retain a fixed gap in the exponent, and
only then absorb finitely many positive contributions. It avoids a fragile
sequence of unnecessarily sharp numerical thresholds. Our implementation
keeps the relevant scale comparisons and the final assembly in separate
modules.

\section{Other Erd\H{o}s problems: concrete connections and limitations}\label{sec:other}
The links below are proposed uses of the machinery, not new results on the
listed problems. The statements are taken from the original problem literature
and the modern problem catalogue, with contemporary comparisons to
Tao--Ter\"av\"ainen and Lau. A common arithmetic vocabulary does not make an
existence theorem a consequence of an upper-counting bound.

\subsection{Problem \#679: smaller backward prime-factor barriers}
Problem \#679 asks, for each $\varepsilon>0$, whether infinitely many $n$ satisfy
\[
 \omega(n-k)<(1+\varepsilon)\frac{\log k}{\log\log k}
\]
for all $k<n$ sufficiently large depending only on $\varepsilon$
\citep{Erdos1979Acta,Bloom679}. The quantifier on the initial cutoff is important:
it must be independent of $n$.

This is a close structural match to the backward-window bookkeeping. A possible
first application is a finite necessary-condition count over shifts
$K_\varepsilon\le k\le H$, with the aggregate barrier replacing $Q_H$ and the
mandatory debit recomputed for the shorter shifted window. The weighted
coefficient and CRT machinery would then be available. Such an application
could quantify the scarcity of integers satisfying a prescribed finite set of
barriers. It would not, by itself, settle whether infinitely many full barrier
solutions exist, nor may one let $K_\varepsilon$ grow with $n$ without changing
the question. Lau discusses sharper barriers in connection with a conditional
Cram\'er-type model; those conditional conclusions are not inputs to our proof
\citep{Lau2026}.

\subsection{Problem \#826: linear forward divisor barriers}
Problem \#826 asks whether infinitely many $n$ have
$\tau(n+k)\ll k$ for every $k\ge1$, with a uniform implied constant
\citep{Bloom826}. Lau obtains a polynomial barrier
$\tau(n+k)\ll k^C$ for an absolute $C$, not the requested linear one
\citep{Lau2026}.

A forward window admits the same local congruence counting after reversing the
shift convention. For a fixed proposed divisor-barrier constant, the
logarithmic-budget conversion could yield finite upper bounds for its
survivors, with constants and the length of the window kept explicit.
The limitation is decisive: $2^{\omega(m)}\le\tau(m)$ gives a necessary
condition for small $\tau$, not a sufficient one. Constructing many integers
with small divisor counts also requires control of prime multiplicities and
cofactors. The current sieve does not provide that construction.

\subsection{Problem \#413: the barrier \texorpdfstring{$\omega(n-k)\le k$}{omega(n-k) <= k}}
Problem \#413 concerns infinitely many integers obeying backward prime-factor
barriers of the form $\omega(n-k)\le k$ for $1\le k<n$
\citep{Bloom413,TaoTeravainen2026}. The same finite incidence-counting and
congruence modules could formalize a relaxation of these constraints.
But summing this barrier gives a budget of order $H^2$, rather than the
$H\log H$ budget used in this paper. Thus neither the exponent $a$ nor the
mass--budget gap transfers automatically. A useful project here would be a
formal, parameter-dependent finite counting theorem for that exact barrier,
followed by an independent assessment of whether it says anything substantive
about the infinite existence question.

\subsection{Problem \#248: a solved benchmark for reusable formalization}
Problem \#248 asks for infinitely many $n$ with a uniform linear forward bound
on the number of prime factors. Tao--Ter\"av\"ainen prove
$\omega(n+k)\le\Omega(n+k)\ll k$ for every positive $k$
\citep{Bloom248,TaoTeravainen2026}; Lau subsequently obtains a logarithmic
barrier for $k\ge2$ \citep{Lau2026}.

Here the prospective benefit is formalization, not another claimed solution.
The finite weighted algebra, local root counting, and exponential-error
interfaces could serve as independently usable components of a formalized
proof or companion quantitative estimate. The existence arguments in those
papers need additional probabilistic and sieve ingredients. Our deterministic
upper-moment estimate cannot replace those ingredients merely because both
arguments discuss prime factors of nearby integers.

\subsection{A disciplined transfer criterion}
For each proposed application, three tasks should precede a large new
formalization: identify the exact arithmetic barrier and its quantifiers;
compute the local root densities and compulsory debit; and determine whether
the desired conclusion is an upper bound, an existence result, or a uniform
all-shifts assertion. Only then is reuse more than a similarity of notation.
The present work is strongest as a library for finite counting and audited
analytic assembly. It does not furnish a general mechanism for resolving all
prime-factor problems.

\section{Conclusion}\label{sec:conclusion}
The divisor restrictions in Erd\H{o}s \#647 impose a logarithmic bound on the
total number of distinct prime-factor occurrences in a growing backward
window. Removing the unavoidable small-prime occurrences leaves a signed
budget that can be compared with a precisely evaluated selected-prime mass.
A finite weighted Bonferroni argument, exact CRT root counts, and an injective
encoding of all retained prime subsets control the averaging error without an
extra factor depending on the truncation order.

At the exponent $a=\log2/(1+\log2)$, the leading terms cancel. Mertens error
cancellation and the retained linear factorial term leave a fixed positive
mass--budget gap. The complete amplified principal term, factorial tail,
arithmetic remainder, and exceptional window then give
Theorem~\ref{thm:majorant}. Final ratio absorption proves
Theorem~\ref{thm:main} for every fixed $0<c<1499/10^6$, including $1/1000$.

The audited formalization and the human-readable proof have complementary
roles. Lean fixes exact hypotheses, signs, boundary cases, and dependencies;
the exposition identifies the mathematical mechanism and its limitations.
The AI--operator workflow produced the development, but its claims rest on
the statements and proof artifacts rather than on the model's authority.
The reusable finite and elementary library is separate from the analytic
package, and the latter imports only the audited PNT+ slice it needs.

What remains is not a missing step in this endpoint deduction, but a different
question: can one exclude every candidate beyond $24$, or construct one?
The bound proved here still allows an infinite sparse set. External comparison
of novelty, mathematical review, broader independent proof replay, and any
extension to related problems should therefore be pursued as distinct tasks,
not conflated with a resolution of \#647.

\section*{Acknowledgments and contributions}
\addcontentsline{toc}{section}{Acknowledgments and contributions}
The project uses the work of the Lean, Mathlib, and PNT+ communities, and the
public problem curation of Thomas F. Bloom. Existing mathematical and formal
sources are credited in the bibliography and the repository. GPT 6 Astra is
the AI author attribution requested for this paper. Justin V. Wick operated
the interactive workflow, executed the verification runs, supplied the returned
artifacts, and managed the project. No institutional affiliation or endorsement
by the cited communities is implied.

\appendix
\section{Technical details and boundary cases}\label{app:technical}
\subsection{The finite Bonferroni recurrence}\label{app:bonferroni}
This supplies a direct algebraic proof of \eqref{eq:bonferroni}. For a finite
list $w$, write $V(w)=\prod_i(1-w_i)$ and $E_J(w)=P_J(w)-V(w)$. Inserting a
new weight $u$ gives
\[
 e_{j+1}(u::w)=e_{j+1}(w)+u e_j(w),\qquad
 P_{j+1}(u::w)=P_{j+1}(w)-uP_j(w).
\]
It follows that
\[
 E_{j+1}(u::w)=E_{j+1}(w)-uE_j(w).
\]
The invariant $(-1)^j E_j(w)\ge0$ is proved by induction on the list. At
$j=0$, $E_0=1-V\ge0$ because $0\le V\le1$. At $j+1$, multiplying the
recurrence by $(-1)^{j+1}$ expresses it as the sum of the two nonnegative
quantities $(-1)^{j+1}E_{j+1}(w)$ and $u(-1)^jE_j(w)$.
For even $J$ this proves $P_J\ge V$. The next truncation satisfies
$P_{J+1}=P_J-e_{J+1}\le V$, giving the upper half of
\eqref{eq:bonferroni}. This also covers the empty list and orders above its
length.

\subsection{The elementary logarithm constants}\label{app:constants}
The numerical arguments need only coarse certified bounds
\[
 0.693<\log2<0.7.
\]
These follow from the explicit real-logarithm bounds in the pinned Mathlib
sources \citep{MathlibPin}. Since $25/2>8$,
\[
 K>3\log2+\frac1{\log2}-2
   >3\frac{693}{1000}+\frac{10}{7}-2
   =\frac{10553}{7000}>\frac32.
\]
Similarly $20>16$ implies $\log20>4\log2>2772/1000$, whence
\[
 \frac{\log20-1}{50}-\frac1{500}
 >\frac{1.772}{50}-\frac1{500}
 =\frac{209}{6250}>\frac1{30}.
\]
All decimals in these displayed comparisons denote exact rational numbers.
No floating-point numerical evaluation is a proof dependency. Sharpening these
inequalities is unnecessary for the stated endpoint.

\subsection{Boundary cases recorded in the formal interfaces}
\begin{center}
\begin{tabularx}{\textwidth}{@{}p{0.25\textwidth}X@{}}
\toprule
Case & Treatment \\
\midrule
Negative interval translation & The finite moment uses integer intervals and Euclidean quotient counts; no positivity is imposed on the translation. \\
Zeros of $F_H(n)$ & Divisibility and CRT root identities apply to all integer inputs, including zero values of the polynomial. \\
Empty prime set or subset & Empty products equal one; the empty subset has modulus one and exact count $X$. \\
$J=0$ or $J>|\Pset|$ & The finite Bonferroni and coefficient conventions cover these cases. Endpoint parameters eventually have $J>0$. \\
Negative corrected budget & A separate exclusion gives $C(X)\le H$; the positive-weight estimates are not multiplied by a negative budget. \\
Coefficient at the boundary & Only $c<\rate$ is asserted. Equation~\eqref{eq:boundary-ratio} shows the limitation of the current majorant. \\
A translated endpoint & The finite moment is uniform in translation. No separate all-translations endpoint theorem is claimed as an exported result here. \\
\bottomrule
\end{tabularx}
\end{center}

\section{Formalization map and verification provenance}\label{app:evidence}
\subsection{Theorem-to-source correspondence}\label{app:map}
The following map identifies representative exported or audited declarations;
it is not a claim that every sentence of the exposition is a separate Lean
lemma. In the table, the prefixes are
$\mathsf{S}$ for \smallcode{Erdos647Sieve},
$\mathsf{E}$ for \smallcode{Erdos647Sieve.Elementary},
$\mathsf{A}$ for \smallcode{Erdos647Sieve.Analytic}, and
$\mathsf{P}$ for \smallcode{Erdos647Sieve.Endpoint}.
The complete source snapshot is in the accompanying verification artifact
\citep{SieveArtifact}.

\begingroup\small
\renewcommand{\arraystretch}{1.18}
\begin{longtable}{@{}>{\raggedright\arraybackslash}p{0.29\textwidth}>{\raggedright\arraybackslash}p{0.66\textwidth}@{}}
\toprule
Mathematical component & Representative Lean declaration \\
\midrule
\endfirsthead
\toprule
Mathematical component & Representative Lean declaration \\
\midrule
\endhead
\bottomrule
\endfoot
Candidate and finite claims & $\mathsf{S}$.\leanname{Candidate};
 $\mathsf{S}$.\leanname{FiniteMomentClaim};
 $\mathsf{S}$.\leanname{FiniteCountingClaim} \\
Mandatory debit, Section~\ref{sec:budget} &
 $\mathsf{S}$.\leanname{budgetDebit} \\
Weighted inclusion--exclusion, Section~\ref{sec:moment} &
 $\mathsf{S}$.\leanname{Bonferroni.even_sum_bounds} \\
Finite moment and transfer, Sections~\ref{sec:moment}--\ref{sec:transfer} &
 $\mathsf{S}$.\leanname{finiteMoment};
 $\mathsf{S}$.\leanname{finiteCounting} \\
Reciprocal-prime lower bound &
 $\mathsf{E}$.\leanname{prime_reciprocal_lower} \\
Fixed-constant debit &
 $\mathsf{E}$.\leanname{smallPrimeDebit_lower} \\
Factorial lower bound and shifted limit &
 $\mathsf{E}$.\leanname{factorial_pow_lower};
 $\mathsf{E}$.\leanname{shiftedFactorialLog_residual_tendsto_zero} \\
Budget estimate, Proposition~\ref{prop:budget-upper} &
 $\mathsf{E}$.\leanname{correctedBudget_upper_eventually} \\
Mertens convergence &
 $\mathsf{A}$.\leanname{reciprocalPrimeLimit_exists} \\
Two-cutoff error cancellation &
 $\mathsf{A}$.\leanname{selectedPrimeSum_difference_tendsto_zero} \\
Actual prime-mass expansion &
 $\mathsf{P}$.\leanname{endpoint_primeMass_expansion_tendsto_zero} \\
Positive mass--budget gap &
 $\mathsf{P}$.\leanname{eventually_primeMass_sub_correctedBudget_ge} \\
Mass and budget sizes &
 $\mathsf{P}$.\leanname{eventually_endpoint_mass_budget_bounds} \\
Principal exponent &
 $\mathsf{P}$.\leanname{eventually_amplification_exponents_le} \\
Amplified factorial tail &
 $\mathsf{P}$.\leanname{eventually_amplified_moment_tail_le} \\
All-sign counting bound &
 $\mathsf{P}$.\leanname{eventually_candidateCount_le_amplified_errors} \\
Final absorption &
 $\mathsf{P}$.\leanname{eventually_amplifiedErrorMajorant_le_endpointRHSWith} \\
General coefficient theorem &
 $\mathsf{P}$.\leanname{endpointBound_of_lt_principal_rate} \\
Explicit and historical corollaries &
 $\mathsf{P}$.\leanname{endpointBound_one_div_thousand};
 $\mathsf{P}$.\leanname{positiveEndpoint};
 $\mathsf{S}$.\leanname{endpoint} \\
\end{longtable}
\endgroup

The final assembly is intentionally short. In outline, it combines the
all-sign majorant with its absorption into the desired endpoint expression:
\begin{quote}\small\ttfamily
 theorem endpointBound\_of\_lt\_principal\_rate\par
\hspace*{1em}(c : Real) (hc : c < (1499 / 1000000 : Real)) :\par
\hspace*{1em}EndpointBound c
\end{quote}
The printed type is not the proof; the supplied source contains the proof
body, and the final audit checks both the expanded endpoint definition and the
transitive axioms. In particular $c$ is not an assumed bound on a finite
majorant: that bound has already been proved in the imported absorption module.

\subsection{Exact snapshot and trust record}
The development URL is \url{https://github.com/JustinWick/erdos647}.
The paper's precise basis is the operator-supplied v48 archive identified below,
rather than an assertion that a mutable repository branch always contains the
same bytes. The recorded UTC run spans
\texttt{2026-09-09 01:40:11}--\texttt{01:45:24}.

\begin{center}
\begin{tabularx}{\textwidth}{@{}p{0.38\textwidth}X@{}}
\toprule
Recorded item & Value \\
\midrule
Lean toolchain & \smallcode{leanprover/lean4:v4.32.2} \\
Platform in compiler report & \smallcode{x86\_64-unknown-linux-gnu} \\
Configured proof gates & 30 passed; none failed or blocked \\
Distinct audited declarations & 235 \\
Exact-type/definition examples & 220 in 32 audit files \\
Engineering tests & 317 passed \\
Executed commands & 70, all with zero exit status \\
Build and audit diagnostics & No warnings or errors found in the reviewed logs \\
Archive manifest & 296 entries, all matching their recorded hashes \\
Source snapshots & 179 checked inputs before and after; unchanged \\
Retained PNT+ source closure & 14 modules; admission scan passed \\
Axiom policy & Exactly the allowed set in \eqref{eq:axioms}; no unauthorized axiom reported \\
\bottomrule
\end{tabularx}
\end{center}

\noindent\textbf{Version identifiers.}
\begin{quote}\small
Lean compiler commit:\par
\nolinkurl{f3b06c705e6c85f5314019d5d3baab0fec5b580c}\par\medskip
Mathlib revision:\par
\nolinkurl{905b95818eb32af7874a58b427f50c1711a5e96c}\par\medskip
PNT+ revision:\par
\nolinkurl{a5154676af9aa3095150ee410cdda80555aa0642}
\end{quote}

\noindent\textbf{Artifact and source hashes.}
\begin{quote}\small
Run archive:\par
\nolinkurl{erdos647_repo_v48_results_20260909T014011Z_dd165b9f.zip}\par
SHA-256:\par
\nolinkurl{7043421f1ec2ca5c85963e84a1a670096b23a2688fb16932a42b738c1343561e}\par\medskip
\nolinkurl{Erdos647Sieve/Specification.lean}:\par
\nolinkurl{39d7c9e69faba89c27d738c6bf37f8b363197e2e5942281c5aedbbba7374c574}\par\medskip
\nolinkurl{research/Erdos647Research/Endpoint/Main.lean}:\par
\nolinkurl{91700a3fe721917602595cbd80ba1a515ada11abc679106b37f4167426dec0b4}
\end{quote}
The review that produced these counts verifies the supplied archive and its
logs; it is not another Lean execution. The final endpoint had already passed
in the v47 proof checkpoint. The v48 record additionally checks the promoted
Mathlib-only elementary modules and the public analytic and endpoint facades.
Thus the packaging changes are not confused with a new mathematical theorem.

\subsection{Reproducibility and the public interface}
The repository contains pinned toolchain and Lake manifests, ordinary core
build/test targets, and separate research targets. The top-level runner manages
incremental checking, setup when explicitly requested, and diagnostic archives.
A mathematical consumer can instead import the public facades without knowing
the historical overlay sequence. The core and analytic packages maintain
separate dependency boundaries; only the latter needs PNT+.

The companion paper-source bundle includes this LaTeX source, bibliography,
a PDF build script, the evidence-review report, and the original v48 diagnostic
archive. It is sufficient to identify the proof snapshot described here.
Public deployment, independent kernel replay, and peer review are not inferred
from the presence of a CI configuration or from file hashes.

\section{Notation and proof dependencies}\label{app:notation}
\begin{center}
\begin{tabularx}{\textwidth}{@{}p{0.19\textwidth}X@{}}
\toprule
Symbol & Meaning \\
\midrule
$\tau,\omega,\Omega$ & Divisor count; number of distinct prime factors; number with multiplicity \\
$C(X)$ & Number of actual full-condition candidates $24<n\le X$ \\
$F_H(n)$ & $\prod_{k=1}^H(n-k)$, the window polynomial \\
$\Pset(H,y)$ & Primes $H<p\le y$ \\
$T_{H,y}(n)$ & Number of selected primes dividing $F_H(n)$ \\
$Q_H,D_H,B_H$ & Logarithmic budget, mandatory small-prime debit, signed difference \\
$\lambda$ & Prime mass $H\sum_{H<p\le y}1/p$ \\
$Q,L,U,S$ & $\log X$, $\log\log X$, $Q^a$, $U/L$ \\
$H,J,y$ & $\floor{Q^a}$, $2\ceil{H/100}$, $X^{1/(4J)}$ \\
$t,z,\eta,\mu$ & $1/(1000L)$, $1-t$, $-\log z$, $t\lambda$ \\
$K$ & $\log(25/2)+1/\log2-2$, the available gap constant \\
$\rate$ & $1499/10^6$, the upper threshold of the proved open coefficient range \\
$c$ & A fixed positive coefficient strictly below $\rate$ \\
\bottomrule
\end{tabularx}
\end{center}

The logical dependencies may be read as two converging chains. The finite
chain is: weighted Bonferroni algebra and CRT progression counts imply the
finite moment; that moment together with the signed divisor budget implies the
finite counting inequality. The analytic chain is: the PNT/integrability input
implies Mertens convergence; actual cutoff asymptotics and the elementary
factorial/debit estimates imply the mass--budget gap. The gap controls the
amplification of the finite inequality, and the scale limits absorb all
remaining contributions. No arrow in these chains stands for an unproved
premise in the final Lean theorem.

\addcontentsline{toc}{section}{References}
\bibliographystyle{plainnat}
\bibliography{references}
\end{document}
